% !TEX program = pdflatex
\documentclass[10pt]{article}

\usepackage[utf8]{inputenc}
\usepackage[T1]{fontenc}
\usepackage[english]{babel}
\usepackage{lmodern}
\usepackage{microtype}
\usepackage[margin=0.72in]{geometry}
\usepackage{amsmath,amssymb,amsthm,mathtools}
\usepackage{booktabs}
\usepackage{xcolor}
\usepackage{enumitem}
\usepackage[most]{tcolorbox}
\usepackage{fancyhdr}
\usepackage[hidelinks]{hyperref}

\definecolor{CourseBlue}{HTML}{1A3C68}
\definecolor{CourseGreen}{HTML}{216E39}
\definecolor{CourseGold}{HTML}{B8860B}
\definecolor{SoftBlue}{HTML}{F2F6FA}
\definecolor{SoftGreen}{HTML}{F1F8F3}
\definecolor{SoftGold}{HTML}{FFF8E8}

\setlength{\parindent}{0pt}
\setlength{\parskip}{0.32em}
\setlist[itemize]{leftmargin=1.5em,itemsep=0.12em,topsep=0.2em}
\setlist[enumerate]{leftmargin=1.7em,itemsep=0.12em,topsep=0.2em}
\allowdisplaybreaks

\pagestyle{fancy}
\fancyhf{}
\lhead{Stochastic Calculus and Option Pricing}
\rhead{Lecture 1 Notes}
\cfoot{\thepage}
\renewcommand{\headrulewidth}{0.3pt}

\newcommand{\E}{\mathbb E}
\newcommand{\Pbb}{\mathbb P}
\newcommand{\Var}{\operatorname{Var}}
\newcommand{\Cov}{\operatorname{Cov}}
\newcommand{\SD}{\operatorname{SD}}
\newcommand{\R}{\mathbb R}
\newcommand{\1}{\mathbf 1}
\newcommand{\cF}{\mathcal F}
\newcommand{\cG}{\mathcal G}
\newcommand{\cH}{\mathcal H}
\newcommand{\dd}{\,\mathrm d}

\theoremstyle{definition}
\newtheorem{definition}{Definition}[section]
\newtheorem{proposition}[definition]{Proposition}
\newtheorem{theorem}[definition]{Theorem}
\newtheorem{remark}[definition]{Remark}

\newtcolorbox{boardexample}[1]{
  enhanced,breakable,colback=SoftBlue,colframe=CourseBlue,
  boxrule=0.55pt,arc=1mm,left=1.5mm,right=1.5mm,top=1mm,bottom=1mm,
  before skip=0.55em,after skip=0.55em,
  title={Example --- #1},fonttitle=\bfseries}
\newtcolorbox{studentexercise}[1]{
  enhanced,breakable,colback=SoftGold,colframe=CourseGold,
  boxrule=0.55pt,arc=1mm,left=1.5mm,right=1.5mm,top=1mm,bottom=1mm,
  before skip=0.55em,after skip=0.35em,
  title={Exercise --- #1},fonttitle=\bfseries}
\newtcolorbox{shortsolution}[1][]{
  enhanced,breakable,colback=SoftGreen,colframe=CourseGreen,
  boxrule=0.55pt,arc=1mm,left=1.5mm,right=1.5mm,top=1mm,bottom=1mm,
  before skip=0.25em,after skip=0.6em,
  title={Solution#1},fonttitle=\bfseries}
\newtcolorbox{keyidea}[1]{
  enhanced,breakable,colback=white,colframe=CourseGreen,
  boxrule=0.7pt,arc=1mm,left=1.5mm,right=1.5mm,top=1mm,bottom=1mm,
  before skip=0.5em,after skip=0.5em,
  title={#1},fonttitle=\bfseries}

\title{\textbf{Option Pricing: Motivation and First Mathematical Tools}\\
\large Lecture 1 Notes --- Class Support and Student Summary}
\author{}
\date{}

\begin{document}
\maketitle
\vspace{-1.4em}


% ============================================================
\section{Motivation, Replication, and Pricing}

\subsection{Contracts, payoffs, and discounting}

A \emph{financial contract} specifies future cash flows. A derivative links
those cash flows to an observable market quantity. For a European call with
strike $K$ and maturity $T$,
\[
X=(S_T-K)^+,
\qquad (x)^+=\max(x,0).
\]
The long position receives $X$ at $T$ and the short position owes $X$.

If the continuously compounded rate is constant and equal to $r$, the bank
account is $B_t=e^{rt}$. A certain payment $C$ at time $T$ has present value
$Ce^{-rT}$.

\subsection{The casino example: expectation is not a price}

\begin{boardexample}{historical expectation versus replication}
A contract pays $\$100$ after red and $\$0$ after black. Historically,
$\Pbb(\text{red})=0.99$, so its physical expectation is $\$99$.

A $\$50$ casino bet pays $\$100$ after red and $\$0$ after black. It reproduces
the contract state by state. Therefore the replication cost is $\$50$, not
$\$99$.
\end{boardexample}

\begin{proposition}[Law of one price]
If two portfolios have the same payoff in every state, then they must have the
same price in an arbitrage-free market.
\end{proposition}

\begin{proof}
If portfolio $A$ were more expensive than $B$, sell $A$ and buy $B$. The initial
cash flow is positive and the terminal cash flow is zero in every state. This
is an arbitrage. Interchange $A$ and $B$ for the opposite price inequality.
\end{proof}

The casino example isolates the central principle:
\[
\boxed{\text{price of a replicable claim}=\text{cost of replication}.}
\]
Historical probabilities describe frequencies; they do not override a
state-by-state replication argument.

\subsection{A one-period financial market}

Let a stock satisfy $S_0>0$ and
\[
S_1=\begin{cases}S_u,&\text{up},\\S_d,&\text{down},\end{cases}
\qquad S_u>S_d,
\]
and let the risk-free gross return be $R>0$. A claim pays $H_u$ or $H_d$.

\begin{proposition}[Replication and the unique one-period price]
The portfolio with $\Delta$ shares and maturity cash $B_1$ replicates $H$ when
\[
\Delta=\frac{H_u-H_d}{S_u-S_d},
\qquad
B_1=H_d-\Delta S_d.
\]
Its no-arbitrage price is
\[
V_0=\Delta S_0+\frac{B_1}{R}.
\]
\end{proposition}

\begin{proof}
The two replication equations are
$\Delta S_u+B_1=H_u$ and $\Delta S_d+B_1=H_d$. Subtract them to obtain
$\Delta$, then substitute to obtain $B_1$. The law of one price identifies
$V_0$ with the initial cost of this portfolio.
\end{proof}

Define
\[
q=\frac{RS_0-S_d}{S_u-S_d}.
\]
The no-arbitrage inequalities $S_d<RS_0<S_u$ are equivalent to $0<q<1$.
Moreover,
\[
S_0=\frac1R\bigl(qS_u+(1-q)S_d\bigr),
\qquad
V_0=\frac1R\bigl(qH_u+(1-q)H_d\bigr).
\]
Thus $q$ is the \emph{risk-neutral up probability}. It is determined by traded
prices, not by the physical frequency of the up state.

\begin{boardexample}{replicate a call}
Assume $R=1$, $S_0=50$, $S_u=60$, $S_d=40$, and $K=50$. The call pays
$(10,0)$. Hence
\[
\Delta=\frac{10-0}{60-40}=\frac12,
\qquad
B_1=0-\frac12(40)=-20,
\qquad
V_0=\frac12(50)-20=5.
\]
Also $q=(50-40)/(60-40)=1/2$, so
$V_0=\E^q[H]=\frac12(10)+\frac12(0)=5$.
\end{boardexample}

\begin{studentexercise}{a similar one-period call}
Assume $R=1$, $S_0=100$, $S_u=125$, $S_d=75$, and $K=100$.
Find the replicating portfolio, price, and risk-neutral probability.
\end{studentexercise}
\begin{shortsolution}
The call payoff is $25$ in the up state and $0$ in the down state. Therefore
the replicating portfolio must satisfy
\[
125\Delta+B_1=25,
\qquad
75\Delta+B_1=0.
\]
Subtracting the second equation from the first gives
\[
50\Delta=25,
\qquad \Delta=\frac12.
\]
Substitution into the down-state equation gives
$B_1=-75/2=-37.5$. Since $R=1$, the time-zero cost is
\[
V_0=\Delta S_0+\frac{B_1}{R}
=\frac12(100)-37.5=12.5.
\]
Finally, the risk-neutral probability is determined by the stock price:
\[
100=q(125)+(1-q)75,
\qquad q=\frac12.
\]
This gives the same claim value,
$V_0=\E^q[H]=\frac12(25)+\frac12(0)=12.5$.
\end{shortsolution}

% ============================================================
\section{Probability Toolbox}

\subsection{Random variables, distributions, CDFs, and densities}

\begin{definition}[Random variable and distribution]
A random variable is a measurable map $X:(\Omega,\cF)\to(\R,\mathcal B(\R))$.
Its distribution is the probability measure
\[
\Pbb_X(A)=\Pbb(X\in A),\qquad A\in\mathcal B(\R).
\]
\end{definition}

\begin{definition}[CDF and density]
The cumulative distribution function is
\[
F_X(x)=\Pbb(X\le x).
\]
It is nondecreasing, right-continuous, and satisfies
$F_X(-\infty)=0$, $F_X(+\infty)=1$. If a nonnegative function $f_X$ satisfies
\[
F_X(x)=\int_{-\infty}^x f_X(u)\dd u,
\]
then $f_X$ is a density of $X$. At differentiability points, $f_X=F_X'$.
\end{definition}

For $a<b$,
\[
\Pbb(a<X\le b)=F_X(b)-F_X(a).
\]
A distribution always exists; a density need not exist.

\begin{boardexample}{a payoff is a random variable}
Suppose
\[
\begin{array}{c|ccc}
S_T&80&100&130\\ \hline
\Pbb&0.20&0.50&0.30
\end{array}
\]
and $X=(S_T-100)^+$. Then $X$ takes values $0,0,30$, so
\[
\Pbb(X=0)=0.70,
\qquad
\Pbb(X=30)=0.30.
\]
The payoff rule maps states into numbers; the probabilities of those numbers
form its distribution.
\end{boardexample}

\begin{boardexample}{recover the exponential CDF}
If $X\sim\operatorname{Exp}(\lambda)$, then
$f_X(x)=\lambda e^{-\lambda x}\1_{\{x\ge0\}}$. Therefore
\[
F_X(x)=\begin{cases}
0,&x<0,\\
1-e^{-\lambda x},&x\ge0.
\end{cases}
\]
Consequently, for $0\le a\le b$,
$\Pbb(a<X\le b)=e^{-\lambda a}-e^{-\lambda b}$.
\end{boardexample}

\begin{studentexercise}{CDF of a one-day return}
Let $\Pbb(R=-2\%)=0.40$ and $\Pbb(R=3\%)=0.60$. Write $F_R$ and compute
$\Pbb(R>0)$.
\end{studentexercise}
\begin{shortsolution}
The CDF accumulates the probability mass at and below the point $r$. Before
$-2\%$ no value is possible; at $-2\%$ it jumps by $0.40$, and at $3\%$ it
jumps by the remaining $0.60$. Hence
\[
F_R(r)=\begin{cases}
0,&r<-2\%,\\
0.40,&-2\%\le r<3\%,\\
1,&r\ge3\%,
\end{cases}.
\]
Since $0$ lies between the two possible returns, $F_R(0)=0.40$. Therefore
\[
\Pbb(R>0)=1-F_R(0)=1-0.40=0.60.
\]
Equivalently, the only positive outcome is $R=3\%$, which has probability
$0.60$.
\end{shortsolution}

\subsection{Expectation, variance, and dependence}

\begin{definition}[Expectation and moments]
For integrable $X$,
\[
\E[X]=\sum_x x\Pbb(X=x)
\quad\text{or}\quad
\E[X]=\int_{\R}x f_X(x)\dd x.
\]
More generally, whenever it exists,
\[
\E[g(X)]=\sum_x g(x)\Pbb(X=x)
\quad\text{or}\quad
\E[g(X)]=\int_{\R}g(x)f_X(x)\dd x.
\]
The $k$th moment is $\E[X^k]$. In particular,
\[
\mu=\E[X],
\quad m_2=\E[X^2],
\quad \Var(X)=\E[(X-\mu)^2],
\quad \SD(X)=\sqrt{\Var(X)}.
\]
\end{definition}

\begin{proposition}[Fundamental expectation properties]
For integrable $X,Y$ and constants $a,b$,
\[
\E[aX+bY]=a\E[X]+b\E[Y],
\qquad
X\le Y\Longrightarrow \E[X]\le\E[Y],
\qquad
|\E[X]|\le\E[|X|].
\]
Also $\E[\1_A]=\Pbb(A)$.
\end{proposition}

\begin{proof}
Linearity follows from linearity of sums or integrals. If $X\le Y$, then
$Y-X\ge0$ and $\E[Y-X]\ge0$. Finally,
$-|X|\le X\le|X|$ and monotonicity give the absolute-value inequality.
The indicator identity follows directly from its two possible values.
\end{proof}

\begin{proposition}[Variance and covariance identities]
\[
\Var(X)=\E[X^2]-\E[X]^2.
\]
With
$\Cov(X,Y)=\E[(X-\E[X])(Y-\E[Y])]$,
\[
\Var(aX+bY)=a^2\Var(X)+b^2\Var(Y)+2ab\Cov(X,Y).
\]
If $X$ and $Y$ are independent and square integrable, then $\Cov(X,Y)=0$.
\end{proposition}

\begin{proof}
Expand $(X-\mu)^2=X^2-2\mu X+\mu^2$ and take expectations. Expanding the
square of $a(X-\E[X])+b(Y-\E[Y])$ gives the second identity. Independence
implies $\E[XY]=\E[X]\E[Y]$, hence zero covariance.
\end{proof}

\begin{center}
\renewcommand{\arraystretch}{1.15}
\begin{tabular}{@{}lccc@{}}
\toprule
Law & Parameters & $\E[X]$ & $\Var(X)$\\
\midrule
Bernoulli & $p$ & $p$ & $p(1-p)$\\
Binomial & $n,p$ & $np$ & $np(1-p)$\\
Poisson & $\lambda$ & $\lambda$ & $\lambda$\\
Uniform & $[a,b]$ & $(a+b)/2$ & $(b-a)^2/12$\\
Exponential & rate $\lambda$ & $1/\lambda$ & $1/\lambda^2$\\
Gaussian & $\mu,\sigma^2$ & $\mu$ & $\sigma^2$\\
\bottomrule
\end{tabular}
\end{center}

\begin{boardexample}{moments of a Bernoulli sum}
Let $B_1,\ldots,B_n$ be independent Bernoulli$(p)$ variables and
$N_n=\sum_{k=1}^nB_k$. Since $B_k^2=B_k$,
\[
\E[B_k]=p,
\qquad \Var(B_k)=p-p^2=p(1-p).
\]
Linearity gives $\E[N_n]=np$. Independence removes cross-covariances, so
$\Var(N_n)=np(1-p)$.
\end{boardexample}

\begin{studentexercise}{classical moment calculations}
Compute $\E[X]$, $\E[X^2]$, and $\Var(X)$ for a uniform and an exponential
variable. Identify the odd moments of a standard Gaussian. As an extension,
derive the Poisson variance from $\E[X(X-1)]$.
\end{studentexercise}
\begin{shortsolution}
For $X\sim\operatorname{Unif}[a,b]$, the density is $1/(b-a)$ on $[a,b]$.
Thus
\[
\E[X]=\frac1{b-a}\int_a^b x\dd x=\frac{a+b}{2},
\qquad
\E[X^2]=\frac1{b-a}\int_a^b x^2\dd x
=\frac{a^2+ab+b^2}{3}.
\]
Consequently,
\[
\Var(X)=\E[X^2]-\E[X]^2=\frac{(b-a)^2}{12}.
\]

For $X\sim\operatorname{Exp}(\lambda)$, integration by parts (or the gamma
integral) gives
\[
\E[X]=\int_0^\infty x\lambda e^{-\lambda x}\dd x=\frac1\lambda,
\qquad
\E[X^2]=\int_0^\infty x^2\lambda e^{-\lambda x}\dd x
=\frac2{\lambda^2}.
\]
Hence $\Var(X)=2/\lambda^2-(1/\lambda)^2=1/\lambda^2$. In summary,
\[
\begin{array}{c|ccc}
X&\E[X]&\E[X^2]&\Var(X)\\ \hline
\operatorname{Unif}[a,b]&(a+b)/2&(a^2+ab+b^2)/3&(b-a)^2/12\\
\operatorname{Exp}(\lambda)&1/\lambda&2/\lambda^2&1/\lambda^2
\end{array}
\]
For $Z\sim\mathcal N(0,1)$, the density is even while $z^{2m+1}$ is odd, so
every odd moment vanishes:
$\E[Z^{2m+1}]=0$. Finally, if
$N\sim\operatorname{Poisson}(\lambda)$, then
\begin{align*}
\E[N(N-1)]
&=e^{-\lambda}\sum_{k=2}^\infty k(k-1)\frac{\lambda^k}{k!}\\
&=\lambda^2e^{-\lambda}\sum_{j=0}^\infty\frac{\lambda^j}{j!}
=\lambda^2.
\end{align*}
Since $N^2=N(N-1)+N$ and $\E[N]=\lambda$, we obtain
$\E[N^2]=\lambda^2+\lambda$ and therefore $\Var(N)=\lambda$.
\end{shortsolution}

\subsection{Sample averages, LLN, and CLT}

For observations $X_1,\ldots,X_n$, define
\[
\overline X_n=\frac1n\sum_{i=1}^nX_i,
\qquad
v_n=\frac1n\sum_{i=1}^n(X_i-\overline X_n)^2,
\qquad
s_n^2=\frac{n}{n-1}v_n.
\]
$v_n$ is the empirical variance; $s_n^2$ is its unbiased version under i.i.d.
sampling with finite variance.

\begin{theorem}[Law of Large Numbers]
If $X_1,X_2,\ldots$ are i.i.d. and $\E[|X_1|]<\infty$, then
\[
\overline X_n\longrightarrow\mu=\E[X_1]
\qquad\text{in probability.}
\]
\end{theorem}

\begin{proof}[Easy proof under the stronger assumption $\Var(X_1)=\sigma^2<\infty$]
Independence gives $\E[\overline X_n]=\mu$ and
$\Var(\overline X_n)=\sigma^2/n$. By Chebyshev's inequality,
\[
\Pbb(|\overline X_n-\mu|>\varepsilon)
\le\frac{\sigma^2}{n\varepsilon^2}\longrightarrow0.
\]
\end{proof}

\begin{theorem}[Central Limit Theorem]
If, in addition, $0<\sigma^2=\Var(X_1)<\infty$, then
\[
\frac{\sqrt n(\overline X_n-\mu)}{\sigma}
\Rightarrow\mathcal N(0,1).
\]
\end{theorem}

The LLN says the estimation error vanishes. The CLT says its typical size is
$\sigma/\sqrt n$ and gives an approximate Gaussian shape. This distinction is
fundamental for statistics, simulation error, and averages of financial data.

\begin{boardexample}{frequency of successes}
For independent Bernoulli$(0.60)$ variables,
\[
\E[\overline X_{100}]=0.60,
\qquad
\SD(\overline X_{100})
=\sqrt{\frac{0.60(0.40)}{100}}\approx0.049.
\]
The CLT gives
\[
\Pbb(0.50\le\overline X_{100}\le0.70)
\approx\Pbb(-2.04\le Z\le2.04)\approx0.96.
\]
\end{boardexample}

\begin{studentexercise}{average daily loss}
Let $X_1,\ldots,X_{225}$ be i.i.d. with mean $2$ and standard deviation $3$.
Find the LLN limit, the standard deviation of $\overline X_{225}$, and a CLT
approximation of $\Pbb(\overline X_{225}>2.4)$.
\end{studentexercise}
\begin{shortsolution}
The variables are i.i.d. with finite mean, so the LLN gives
\[
\overline X_n\xrightarrow[n\to\infty]{\Pbb}2.
\]
In particular, $\overline X_{225}$ is centered at
$\E[\overline X_{225}]=2$ and should be close to $2$ with high probability.
Independence also gives
\[
\Var(\overline X_{225})
=\frac{\Var(X_1)}{225}
=\frac9{225}=0.04,
\qquad
\SD(\overline X_{225})=0.2.
\]
For the tail probability, standardize the threshold using this standard
deviation. The CLT yields
\[
\Pbb(\overline X_{225}>2.4)
\approx\Pbb\!\left(Z>\frac{2.4-2}{0.2}\right)
=\Pbb(Z>2)=1-\Phi(2)\approx0.0228.
\]
\end{shortsolution}

\subsection{Characteristic functions and joint distributions}

\begin{definition}[Characteristic function]
\[
\varphi_X(u)=\E[e^{iuX}],\qquad u\in\R.
\]
It always exists, satisfies $|\varphi_X(u)|\le1$, and uniquely determines the
distribution of $X$.
\end{definition}

\begin{proposition}[Independent sums become products]
If $X$ and $Y$ are independent, then
\[
\varphi_{X+Y}(u)=\varphi_X(u)\varphi_Y(u).
\]
\end{proposition}
\begin{proof}
Independence gives
\[
\E[e^{iu(X+Y)}]=\E[e^{iuX}e^{iuY}]
=\E[e^{iuX}]\E[e^{iuY}].
\]
\end{proof}

A joint distribution records probabilities involving several random variables.
Marginals are obtained by summing or integrating out the other variables. The
joint law, unlike the two marginals alone, records dependence.

% ============================================================
\section{Information and Conditional Expectation}

\subsection{Signals, sigma-algebras, and measurability}

\begin{boardexample}{a signal reveals a partition}
Roll a fair die and let $X(\omega)=\omega$. Observe only the parity signal
$Y=\1_{\{X\text{ is even}\}}$. The observation distinguishes
\[
O=\{1,3,5\},
\qquad
E=\{2,4,6\},
\]
but cannot distinguish outcomes inside either set.
\end{boardexample}

\begin{definition}[Sigma-algebra and information generated by a signal]
A collection $\cG$ of events is a sigma-algebra if
$\Omega\in\cG$, it is closed under complements, and it is closed under
countable unions. For a random variable $Y$,
\[
\sigma(Y)=\{Y^{-1}(B):B\in\mathcal B(\R)\}
\]
is the information revealed by observing $Y$.
\end{definition}

For the parity signal,
$\sigma(Y)=\{\varnothing,O,E,\Omega\}$.

\begin{definition}[Measurable random variable]
A random variable $Z$ is $\cG$-measurable if
\[
\{Z\in B\}\in\cG
\qquad\text{for every Borel set }B.
\]
It is then completely determined by the information in $\cG$.
\end{definition}

\begin{proposition}[Finite-partition test]
If $\cG$ is generated by a finite partition, then $Z$ is $\cG$-measurable if
and only if $Z$ is constant on each atom of that partition. Equivalently,
\[
Z\text{ is }\sigma(Y)\text{-measurable}
\quad\Longleftrightarrow\quad
Z=g(Y)
\text{ for some measurable }g.
\]
\end{proposition}

\begin{proof}
If $Z=g(Y)$, equal values of $Y$ imply equal values of $Z$. Conversely, if $Z$
is constant on every level set of $Y$, define $g(y)$ to be that common value
on $\{Y=y\}$. In a finite model this defines a measurable function $g$.
\end{proof}

\begin{boardexample}{test measurability on parity atoms}
For $Z_1=10-\1_E$, $Z_2=X$, and $Z_3=(-1)^X$:
\begin{itemize}
  \item $Z_1$ is measurable: it is constant on $O$ and on $E$;
  \item $Z_2$ is not measurable: parity does not determine the die value;
  \item $Z_3$ is measurable: it is a function of parity.
\end{itemize}
\end{boardexample}

\begin{studentexercise}{a coarse price signal}
Let $S_T\in\{80,100,130\}$ and $Y=\1_{\{S_T\ge100\}}$. List
$\sigma(Y)$. Decide whether $(S_T-100)^+$ and $\1_{\{S_T\ge100\}}$ are
$\sigma(Y)$-measurable.
\end{studentexercise}
\begin{shortsolution}
The signal is $0$ when $S_T=80$ and $1$ when $S_T\in\{100,130\}$. Its two
information atoms are therefore $\{80\}$ and $\{100,130\}$, and
\[
\sigma(Y)=\{\varnothing,\Omega,\{80\},\{100,130\}\}.
\]
A random variable is $\sigma(Y)$-measurable precisely when it is constant on
each atom. The call payoff $(S_T-100)^+$ equals $0$ at $S_T=100$ and $30$ at
$S_T=130$. It is not constant on the second atom, so it is not
$\sigma(Y)$-measurable. By contrast, the digital payoff is $0$ on $\{80\}$
and $1$ on $\{100,130\}$; it is exactly $Y$ and is therefore measurable.
\end{shortsolution}

\subsection{Conditional distributions and conditional expectation}

For events with $\Pbb(Y=y)>0$,
\[
\Pbb(X=x\mid Y=y)
=\frac{\Pbb(X=x,Y=y)}{\Pbb(Y=y)}.
\]
Conditioning replaces the marginal distribution by the distribution relevant
after the signal has been observed.

\begin{definition}[Conditional expectation]
Let $X\in L^1$ and $\cG\subseteq\cF$. A conditional expectation of $X$ given
$\cG$ is a $\cG$-measurable random variable $Z$ such that
\[
\int_A Z\dd\Pbb=\int_A X\dd\Pbb
\qquad\text{for every }A\in\cG.
\]
It is denoted by $\E[X\mid\cG]$. Conditioning on $Y$ means conditioning on
$\sigma(Y)$:
$\E[X\mid Y]=\E[X\mid\sigma(Y)]$.
\end{definition}

\begin{proposition}[Uniqueness]
Conditional expectation is unique up to almost-sure equality.
\end{proposition}
\begin{proof}
If $Z$ and $Z'$ both satisfy the definition, then
$\int_A(Z-Z')\dd\Pbb=0$ for all $A\in\cG$. Since $Z-Z'$ is measurable, choose
$A=\{Z>Z'\}$ and $A=\{Z<Z'\}$ to conclude $Z=Z'$ almost surely.
\end{proof}

\begin{proposition}[Finite-partition formula]
If $\cG$ is generated by atoms $A_1,\ldots,A_m$ with positive probabilities,
then
\[
\E[X\mid\cG]
=\sum_{j=1}^m\frac{\E[X\1_{A_j}]}{\Pbb(A_j)}\1_{A_j}
=\sum_{j=1}^m\E[X\mid A_j]\1_{A_j}.
\]
\end{proposition}
\begin{proof}
The proposed variable is constant on each atom, hence $\cG$-measurable. Its
integral over $A_j$ equals
$\E[X\mid A_j]\Pbb(A_j)=\E[X\1_{A_j}]$. Add these identities for any union of
atoms.
\end{proof}

\begin{boardexample}{prediction from a partial signal}
On four equally likely states, let $X=(0,4,8,12)$ and let $Y=L$ on states
$\{1,2\}$ and $Y=H$ on states $\{3,4\}$. Then
\[
\E[X\mid Y]
=2\1_{\{Y=L\}}+10\1_{\{Y=H\}}.
\]
The conditional prediction is the average inside the observed information set.
\end{boardexample}

\begin{studentexercise}{a payoff after a coarse die signal}
Roll a fair die and let $X=12\1_{\{\omega=6\}}$. Observe only
$L=\{1,2,3,4\}$ or $H=\{5,6\}$. Compute
$\E[X\mid\sigma(L,H)]$ and verify its expectation.
\end{studentexercise}
\begin{shortsolution}
On $L$, the die cannot equal $6$, so $X=0$ throughout that atom and
$\E[X\mid L]=0$. On $H=\{5,6\}$, the two outcomes remain equally likely and
the corresponding payoffs are $0$ and $12$. Therefore
\[
\E[X\mid H]=\frac{0+12}{2}=6.
\]
The conditional expectation is thus the random variable
\[
\E[X\mid\sigma(L,H)]=6\1_H.
\]
To verify the tower property, note that
\[
\E\!\left[\E[X\mid\sigma(L,H)]\right]
=6\Pbb(H)=6\frac26=2.
\]
Directly, $X=12$ only when the die equals $6$, so
$\E[X]=12(1/6)=2$ as well.
\end{shortsolution}

\subsection{Conditioning with a joint density}

Suppose $(X,Y)$ has joint density $f_{X,Y}$. When $f_Y(y)>0$,
\[
f_Y(y)=\int_{\R}f_{X,Y}(x,y)\dd x,
\qquad
f_{X\mid Y}(x\mid y)=\frac{f_{X,Y}(x,y)}{f_Y(y)}.
\]
For integrable $h$,
\[
g(y)=\int_{\R}h(x)f_{X\mid Y}(x\mid y)\dd x
=\E[h(X)\mid Y=y],
\qquad
\E[h(X)\mid Y]=g(Y).
\]

\begin{boardexample}{the upper triangular density}
Let $f_{X,Y}(x,y)=2\1_{\{0<x<y<1\}}$. For $0<y<1$,
\[
f_Y(y)=\int_0^y2\dd x=2y,
\qquad
f_{X\mid Y}(x\mid y)=\frac1y\1_{(0,y)}(x).
\]
Thus $X\mid(Y=y)\sim\operatorname{Unif}(0,y)$ and
\[
\E[X\mid Y=y]=\frac y2,
\qquad
\E[X\mid Y]=\frac Y2.
\]
\end{boardexample}

\begin{studentexercise}{the reversed triangular density}
Let $f_{X,Y}(x,y)=2\1_{\{0<y<x<1\}}$. Find $f_Y$, $f_{X\mid Y}$, and
$\E[X\mid Y]$.
\end{studentexercise}
\begin{shortsolution}
Fix $y\in(0,1)$. The support condition $0<y<x<1$ allows $x$ to range from
$y$ to $1$. Integrating the joint density over this interval gives
\[
f_Y(y)=\int_y^1 2\dd x=2(1-y).
\]
After normalization,
\[
f_{X\mid Y}(x\mid y)
=\frac{2\1_{(y,1)}(x)}{2(1-y)}
=\frac1{1-y}\1_{(y,1)}(x).
\]
Thus $X\mid(Y=y)\sim\operatorname{Unif}(y,1)$. Its conditional mean can be
read from the midpoint of the interval or calculated directly:
\[
\E[X\mid Y=y]
=\frac1{1-y}\int_y^1x\dd x
=\frac{1-y^2}{2(1-y)}
=\frac{1+y}{2}.
\]
Replacing $y$ by the observed random value $Y$ gives
\[
\E[X\mid Y]=\frac{1+Y}{2}.
\]
\end{shortsolution}

\subsection{Gaussian conditioning}

\begin{proposition}[Conditional law of a Gaussian pair]
Suppose $(X,Y)$ is jointly Gaussian, with
\[
\E[X]=\mu_X,\quad \E[Y]=\mu_Y,\quad
\Var(X)=\sigma_X^2,\quad \Var(Y)=\sigma_Y^2>0,
\]
and $\Cov(X,Y)=\rho\sigma_X\sigma_Y$. Then
\[
X\mid(Y=y)\sim\mathcal N\!\left(
\mu_X+\rho\frac{\sigma_X}{\sigma_Y}(y-\mu_Y),
\;\sigma_X^2(1-\rho^2)
\right).
\]
In particular,
\[
\E[X\mid Y]
=\mu_X+\rho\frac{\sigma_X}{\sigma_Y}(Y-\mu_Y),
\qquad
\Var(X\mid Y)=\sigma_X^2(1-\rho^2).
\]
\end{proposition}

\begin{proof}[Short derivation]
Set
$a=\Cov(X,Y)/\Var(Y)$ and
$R=X-\mu_X-a(Y-\mu_Y)$. Then $\Cov(R,Y)=0$. Since $(R,Y)$ is jointly
Gaussian, zero covariance implies independence. Writing
$X=\mu_X+a(Y-\mu_Y)+R$ gives the conditional mean above, while
$\Var(R)=\sigma_X^2(1-\rho^2)$ gives the conditional variance.
\end{proof}

The conditional mean is the best affine prediction from $Y$. The remaining
variance decreases as $|\rho|$ increases and vanishes when $|\rho|=1$.

\subsection{Properties used in dynamic pricing}

\begin{proposition}[Core conditional-expectation properties]
For integrable variables and appropriate sigma-algebras:
\begin{align*}
\E[aX+bY\mid\cG]&=a\E[X\mid\cG]+b\E[Y\mid\cG],\\
X\le Y&\Longrightarrow \E[X\mid\cG]\le\E[Y\mid\cG],\\
Z\text{ $\cG$-measurable}&\Longrightarrow
\E[ZX\mid\cG]=Z\E[X\mid\cG],\\
X\text{ independent of }\cG&\Longrightarrow
\E[X\mid\cG]=\E[X],\\
\cH\subseteq\cG&\Longrightarrow
\E[\E[X\mid\cG]\mid\cH]=\E[X\mid\cH].
\end{align*}
In particular, $\E[\E[X\mid\cG]]=\E[X]$.
\end{proposition}

\begin{proof}[Proof guide]
For linearity, check the defining integral on every $A\in\cG$ and use
uniqueness. Positivity follows by testing the set where the conditional
expectation would be negative. For a simple known factor
$Z=\sum z_j\1_{A_j}$, pull each constant $z_j$ through the defining integrals;
approximation gives the general case. If $X$ is independent of $\cG$, the
constant $\E[X]$ has the required averages. For the tower rule, test both sides
on every $A\in\cH\subseteq\cG$ and use uniqueness again.
\end{proof}

% ============================================================
\section{Processes, Filtrations, and Martingales}

\subsection{Information evolving through time}

\begin{definition}[Stochastic process]
A discrete-time stochastic process $(X_n)_{n\ge0}$ is a sequence of random
variables on the same probability space. For fixed $n$, $X_n$ is a random
variable; for fixed $\omega$, $n\mapsto X_n(\omega)$ is a sample path.
\end{definition}

\begin{definition}[Filtration and adapted process]
A filtration is an increasing family of sigma-algebras
\[
\cF_0\subseteq\cF_1\subseteq\cF_2\subseteq\cdots.
\]
It represents information available through time. A process $(X_n)$ is adapted
if $X_n$ is $\cF_n$-measurable for every $n$.
\end{definition}

Adapted means ``no look-ahead'': the value at time $n$ can use the observed
past and present, but not a future observation.

\begin{studentexercise}{detect look-ahead}
Let $\xi_1,\xi_2\in\{0,1\}$ encode two coin tosses and
$\cF_n=\sigma(\xi_1,\ldots,\xi_n)$. Consider
\[
X_0=0,
\qquad X_1=\xi_1+\xi_2,
\qquad X_2=\xi_1+\xi_2.
\]
Is $(X_n)$ adapted? Give an adapted cumulative-heads process.
\end{studentexercise}
\begin{shortsolution}
At time $1$, the sigma-algebra $\cF_1=\sigma(\xi_1)$ reveals only the first
toss. Two outcomes can therefore have the same value of $\xi_1$ but different
values of $\xi_2$. On those outcomes, $X_1=\xi_1+\xi_2$ is different, so
$X_1$ is not $\cF_1$-measurable. The original process is consequently not
adapted.

The cumulative number of observed heads is instead
\[
\widetilde X_0=0,
\qquad \widetilde X_1=\xi_1,
\qquad \widetilde X_2=\xi_1+\xi_2.
\]
Here $\widetilde X_n$ depends only on tosses observed by time $n$, so it is
$\cF_n$-measurable for each $n$ and the process is adapted.
\end{shortsolution}

\subsection{Random walks and martingales}

Let
\[
S_n=S_0+\sum_{k=1}^n\varepsilon_k,
\]
where the increments are i.i.d. with mean $m$ and variance $\sigma^2$.

\begin{proposition}[Random-walk moments]
\[
\E[S_n]=S_0+nm,
\qquad
\Var(S_n)=n\sigma^2.
\]
\end{proposition}
\begin{proof}
Use linearity for the mean. For the variance, independence removes all
cross-covariances, leaving the sum of $n$ identical variances.
\end{proof}

\begin{definition}[Martingale]
An integrable adapted process $(M_n)$ is a martingale with respect to
$(\cF_n)$ if
\[
\E[M_{n+1}\mid\cF_n]=M_n
\qquad\text{for every }n.
\]
Given all current information, the best prediction of tomorrow's value is
today's value.
\end{definition}

\begin{boardexample}{the centered random walk is a martingale}
Let $M_n=\sum_{k=1}^n\varepsilon_k$, with independent centered increments and
$\cF_n=\sigma(\varepsilon_1,\ldots,\varepsilon_n)$. Since $M_n$ is known at
time $n$ and $\varepsilon_{n+1}$ is independent of $\cF_n$,
\[
\E[M_{n+1}\mid\cF_n]
=\E[M_n+\varepsilon_{n+1}\mid\cF_n]
=M_n+\E[\varepsilon_{n+1}]
=M_n.
\]
\end{boardexample}

\begin{studentexercise}{a second random-walk martingale}
Let $S_n=\sum_{k=1}^n\varepsilon_k$ for independent symmetric $\pm1$
increments. Show that $M_n=S_n^2-n$ is a martingale.
\end{studentexercise}
\begin{shortsolution}
The variable $M_n=S_n^2-n$ is integrable and $\cF_n$-measurable because it is
a function of the first $n$ increments. It remains to verify the martingale
condition. Since $\varepsilon_{n+1}$ is independent of $\cF_n$ and is
symmetric,
\[
\E[\varepsilon_{n+1}\mid\cF_n]=0,
\qquad
\E[\varepsilon_{n+1}^2\mid\cF_n]=1.
\]
Using $S_{n+1}=S_n+\varepsilon_{n+1}$ and pulling the known quantity $S_n$
out of the conditional expectation,
\begin{align*}
\E[M_{n+1}\mid\cF_n]
&=\E[(S_n+\varepsilon_{n+1})^2-(n+1)\mid\cF_n]\\
&=S_n^2+2S_n\E[\varepsilon_{n+1}\mid\cF_n]
  +\E[\varepsilon_{n+1}^2\mid\cF_n]-(n+1)\\
&=S_n^2-n=M_n.
\end{align*}
Thus all three requirements---integrability, adaptation, and the conditional
expectation identity---are satisfied.
\end{shortsolution}

\begin{keyidea}{Fair games and financial pricing}
In a fair casino game, conditional expected future wealth equals current
wealth. Under a risk-neutral measure $\mathbb Q$, discounted traded-asset
prices are martingales. For a replicable payoff $X$,
\[
\frac{V_t}{B_t}
=\E^{\mathbb Q}\!\left[\frac{X}{B_T}\,\middle|\,\cF_t\right].
\]
Replication supplies the economic restriction; conditional expectation and
martingales express it dynamically.
\end{keyidea}

% ============================================================
\section{Short Differential-Calculus Review}

\begin{definition}[Derivative]
For $g:\R\to\R$, if the limit exists,
\[
g'(x)=\lim_{h\to0}\frac{g(x+h)-g(x)}{h}.
\]
Equivalently, $g(x+h)=g(x)+g'(x)h+o(h)$.
\end{definition}

\begin{definition}[Partial derivatives, gradient, and Hessian]
For $f:\R^d\to\R$ and the $i$th coordinate vector $e_i$,
\[
\partial_i f(x)
=\lim_{h\to0}\frac{f(x+he_i)-f(x)}{h},
\qquad
\nabla f(x)=(\partial_1f(x),\ldots,\partial_df(x)).
\]
If second partial derivatives exist, the Hessian is
\[
\nabla^2f(x)
=\left(\frac{\partial^2f}{\partial x_i\partial x_j}(x)\right)_{i,j}.
\]
\end{definition}

\begin{boardexample}{gradient and Hessian}
For $f(x,y)=x^2y+e^y$,
\[
f_x=2xy,
\qquad f_y=x^2+e^y,
\qquad \nabla f=(2xy,x^2+e^y),
\]
and
\[
\nabla^2f(x,y)
=\begin{pmatrix}
2y&2x\\
2x&e^y
\end{pmatrix}.
\]
\end{boardexample}

\begin{proposition}[Multivariable chain rule]
If $x(t)\in\R^d$ and $f$ are differentiable, then
\[
\frac{\dd}{\dd t}f(x(t))=\nabla f(x(t))\cdot x'(t).
\]
For an option value $V=V(t,S)$ and a differentiable price path $S(t)$,
\[
\frac{\dd}{\dd t}V(t,S(t))
=V_t(t,S(t))+V_S(t,S(t))S'(t).
\]
\end{proposition}

$V_S$ measures sensitivity to the underlying price, $V_t$ sensitivity to
calendar time, and $V_{SS}$ curvature with respect to the underlying.

\begin{studentexercise}{apply the chain rule}
For $V(t,s)=e^{-rt}s^2$ and a differentiable path $S(t)$, compute
$\frac{\dd}{\dd t}V(t,S(t))$.
\end{studentexercise}
\begin{shortsolution}
Differentiate first with respect to each argument while holding the other one
fixed:
\[
V_t(t,s)=-re^{-rt}s^2,
\qquad
V_s(t,s)=2e^{-rt}s.
\]
The multivariable chain rule along the path $t\mapsto(t,S(t))$ gives
\begin{align*}
\frac{\dd}{\dd t}V(t,S(t))
&=V_t(t,S(t))+V_s(t,S(t))S'(t)\\
&=-re^{-rt}S(t)^2+2e^{-rt}S(t)S'(t)\\
&=e^{-rt}\bigl(-rS(t)^2+2S(t)S'(t)\bigr).
\end{align*}
\end{shortsolution}



\end{document}
